\lysh{38860}{4}

\begin{alist}
\item Αν $x$ το μήκος των κάθετων πλευρών του αρχικού γκρι τριγώνου, από το Πυθαγόρειο θεώρημα έχουμε ότι:
$x^2 + x^2 = 16^2$ ή $2x^2 = 256$ ή $x^2 = 128$
και επειδή $x > 0$, $x = \sqrt{128} = \sqrt{64 \cdot 2} = 8\sqrt{2}\text{ cm}$.

\item
\begin{rlist}
\item Έχουμε ότι $y^2 + y^2 = x^2$ δηλαδή $2y^2 = x^2$ άρα $y = \dfrac{x}{\sqrt{2}}$.
Όμοια, έχουμε ότι $z = \dfrac{y}{\sqrt{2}}$ και επειδή $y = \dfrac{x}{\sqrt{2}}$, βρίσκουμε $z = \dfrac{x}{\sqrt{2} \cdot \sqrt{2}} = \dfrac{x}{2}$.
\item Από το ερώτημα $\alpha)$ γνωρίζουμε ότι οι κάθετες πλευρές του αρχικού γκρι, οπότε και η υποτείνουσα του αρχικού άσπρου τριγώνου, έχουν μήκος $8\sqrt{2}\text{ cm}$. Άρα, εφαρμόζοντας το Πυθαγόρειο Θεώρημα στο πρώτο άσπρο τρίγωνο, έχουμε:
$2\beta_1^2 = (8\sqrt{2})^2$
δηλαδή
$2\beta_1^2 = 64 \cdot 2$
οπότε
$\beta_1^2 = 64$
και τελικά $\beta_1 = 8$.

Επίσης, από το ερώτημα $\beta)$i. γνωρίζουμε ότι $z = \dfrac{y}{\sqrt{2}}$ άρα $\dfrac{z}{y} = \dfrac{1}{\sqrt{2}}$, δηλαδή ο λόγος των κάθετων πλευρών δύο διαδοχικών άσπρων τριγώνων είναι σταθερός και ίσος με $\dfrac{1}{\sqrt{2}}$.
Άρα, τα μήκη των κάθετων πλευρών των άσπρων τριγώνων είναι όροι γεωμετρικής προόδου με πρώτο όρο $\beta_1 = 8$ και λόγο $\lambda = \dfrac{1}{\sqrt{2}}$. Ο γενικός όρος της προόδου δίνεται από τη σχέση $\beta_\nu = 8 \left(\dfrac{1}{\sqrt{2}}\right)^{\nu-1}$.

\textit{Άλλος τρόπος:}
Στο ερώτημα $\alpha$ είναι $x = 8\sqrt{2}$. Αν η υποτείνουσα του πρώτου προκύπτοντος άσπρου τριγώνου ταυτίζεται με αυτή του αρχικού γκρι, τότε από το Πυθαγόρειο Θεώρημα για τις δύο ίσες κάθετες πλευρές του ($\beta_1$) ισχύει:
\[2\beta_1^2 = (8\sqrt{2})^2 = 128 \Rightarrow \beta_1^2 = 64 \Rightarrow \beta_1 = 8.\]
\end{rlist}

\item Το συνολικό μήκος $S_\nu$ των κάθετων πλευρών των $\nu$ άσπρων τριγώνων δίνεται από το άθροισμα των όρων γεωμετρικής προόδου με $a_1=8$ και $\lambda=\dfrac{1}{2}$:
\[S_\nu = 8 \dfrac{1 - \left(\dfrac{1}{2}\right)^\nu}{1 - \dfrac{1}{2}} = 16 \left(1 - \left(\dfrac{1}{2}\right)^\nu\right).\]
Επειδή $\left(\dfrac{1}{2}\right)^\nu > 0 \Rightarrow 1 - \left(\dfrac{1}{2}\right)^\nu < 1$, προκύπτει ότι $S_\nu < 16$, άρα και το άθροισμα $\beta$ ικανοποιεί $\beta < 16$.
\end{alist}
